% ============================================================
%  Chapter: Human Values and Value Inculcation
%  Class Notes --- Ethics in Relationships, Human Values,
%  Determinants of Ethics (Family, Education, Society,
%  Government, Media), Lessons from Great Leaders
% ============================================================

\chapter{Human Values and Value Inculcation}


\vspace{14pt}

\section{Recap: Technology Ethics, Media Ethics \& Public vs Private Relationships}

\subsection{Technology Ethics}

Technology ethics are the standards we are supposed to follow when we are developing, deploying, or using technology.

\begin{itemize}
\item While developing technology, we should be mindful not to violate people's privacy.
\item We should be transparent about how the technology is developed and how it is used.
\item We should hold ourselves accountable and responsible for the negative consequences of the technology we develop.
\item Other principles to be followed: beneficence, non-maleficence, and fairness.
\end{itemize}

\begin{KeyConcept}
\begin{itemize}
\item Technology ethics = the standards to follow when developing, deploying or using technology.
\item Core commitments: protect privacy, maintain transparency, uphold accountability for negative consequences, and observe beneficence, non-maleficence and fairness.
\end{itemize}
\end{KeyConcept}

\subsection{Media Ethics}

Media ethics are the principles that media houses should uphold. Media is essentially the medium to transmit information; therefore, information should be verified first before it is broadcast.

\begin{itemize}
\item Fairness and accuracy of information is very important in the context of media ethics.
\item Independence is very important because media houses play the role of the fourth pillar of democracy --- they should be able to hold government accountable.
\item Accountability will not work out if media is not really independent and autonomous.
\end{itemize}

\paragraph{Ways in which media's autonomy is threatened}
\begin{itemize}
\item Government at times misuses its authority --- e.g., filing of fake UAPA cases or sedition cases against media.
\item Print media advertisement is misused as a coercive policy (advertising used as leverage over media houses).
\end{itemize}

\begin{CriticalPoint}
\begin{itemize}
\item Media = fourth pillar of democracy; its accountability role collapses without genuine independence and autonomy.
\item Threats to autonomy: fake UAPA / sedition cases, and misuse of government advertising as a lever of control.
\end{itemize}
\end{CriticalPoint}

\subsection{Ethics in Public vs Private Relationships}

These are the standards we are supposed to follow when dealing with public and private relationships. The essence of the two differs fundamentally.

\paragraph{Nature and essence}
\begin{itemize}
\item Private relationship: its essence is essentially about human emotions and love. Enforcement of ethical standards is rather informal.
\item Public relationship: its essence lies in public good / societal good. Societal good and societal interest become more important than our own individual interest.
\end{itemize}

\paragraph{Source of ethical standards}
\begin{itemize}
\item In private relationships, ethical standards are mostly sourced from traditions and customs.
\item Example: festivals and wedding vows. Across all religions, wedding vows exist --- these vows are nothing but traditions influencing our understanding of our ethical obligation towards our partner and spouse.
\end{itemize}

\paragraph{Why higher ethical obligations apply in public life} When we are discharging a public duty or acting as a public servant, the ethical obligations upon us are higher than when we are only private citizens. This is because our actions have huge ramifications on public interest --- the consequences will be huge if you are a public servant and you are not duty-bound.

\begin{ComparisonBox}
\begin{itemize}
\item Private relationship: essence = emotions \& love; standards informal; sourced from traditions \& customs.
\item Public relationship: essence = public/societal good; higher obligations because actions have huge ramifications on public interest.
\end{itemize}
\end{ComparisonBox}

\section{Ethical Standards in Public Relationships (Nolan Committee Principles)}

The Nolan Committee highlighted certain principles that we are supposed to follow in our public relationships. These are largely the Nolan Committee recommendations.

\begin{GovtPolicy}
\begin{itemize}
\item The Nolan Committee laid down principles of public life, including selflessness, integrity, objectivity, accountability, openness (transparency), and honesty.
\item Objectivity, accountability, openness, integrity and honesty will be discussed in detail under `Foundational Values for Civil Services'.
\end{itemize}
\end{GovtPolicy}

\subsection{Selflessness}

When you are in a public relationship, something else becomes more important than our own individuality --- and that is why selflessness is very important, especially in public relationships.

\begin{itemize}
\item Example --- COVID warriors: It is not that they did not have family or friends to love, nor that they did not love their own life; still they discharged their duty because something is more important --- their duty --- and that is expected when you are in public life.
\item Example --- soldiers: they sacrifice their life for the nation because something else is more important than their own life.
\end{itemize}

\subsection{Integrity}

Integrity is basically following moral principles even when no one is watching.

\begin{itemize}
\item Example --- traffic light: If you jump a red light at night you will not be fined because there is no one to catch you and no one to hold you accountable. If you still do not cross, you are being integral --- you are following the traffic light even when no one is watching.
\item Great leaders are integral in that they follow principles even when no one is watching. In public life, it is important that we become integral.
\end{itemize}

\subsection{Dignity and Compassion}

In public life and public relationships, it is important that we treat each other as human beings with a lot of dignity and compassion.

\begin{itemize}
\item At times our biases enter the picture, and because of those biases class differences operate. All such practices are basically unethical, because every human being deserves dignity.
\item Compassion means feeling others' pain (as felt during Corona). Examples of compassion: Mother Teresa, Swami Vivekananda, Gandhiji.
\end{itemize}

\subsection{Leadership}

In public life it is expected of us that we take leadership --- we should take initiative and not just leave things to others.

\begin{QuickRevision}
\begin{itemize}
\item Public-life ethical standards (Nolan): Selflessness (COVID warriors, soldiers), Integrity (moral principles even when no one is watching), Dignity \& Compassion (every human deserves dignity; compassion = feeling others' pain), Leadership (take initiative).
\item Also to be covered later: objectivity, accountability, openness/transparency, honesty.
\end{itemize}
\end{QuickRevision}

\section{Ethical Standards in Private Relationships}

\subsection{Love and Fidelity}

\begin{itemize}
\item Love: it is expected that there is an element of love in private relationships --- enjoying time with family and friends.
\item Fidelity (loyalty): fidelity becomes an expected ethical standard in private relationships. We judge people on the basis of whether there is loyalty or not.
\end{itemize}

\subsection{Responsibility}

In a family, all of us have responsibilities, and these responsibilities come with our role.

\begin{itemize}
\item A mother is expected to be responsible; if she neglects her child (e.g., when the child meets an accident), we call her irresponsible because responsibility is expected.
\item A father who is an alcoholic and mistreats his wife is called irresponsible because responsibility is expected. Ethics tells you what is expected of you.
\item Fundamental Duties also speak of responsibility: parents have a responsibility towards their children.
\item Children's responsibility towards parents: a viral video showed a boy abandoning his elderly parents at a location while the parents clung to his feet begging him not to leave --- the boy is being irresponsible, because responsibility is expected of children too (you must take care of your elderly parents).
\item An elder sibling has responsibility towards younger siblings. Responsibility is an important part of family and relationships, and it comes with your role.
\end{itemize}

\begin{CriticalPoint}
\begin{itemize}
\item Because official/public duty and private responsibility can both make demands on you, they sometimes clash --- this is the source of ethical dilemmas in case studies.
\item When solving a case study, do NOT pick only one side. Find the golden mean / middle path that addresses BOTH moral alternatives. Going to only one side is not a good answer.
\end{itemize}
\end{CriticalPoint}

\begin{ExampleBox}
\begin{itemize}
\item Sunil case (recurring case study): Sunil is a young, dynamic SP/IPS officer posted in a district infamous for illegal sand mining. He starts taking action against the sand mafia. The mafia threatens to kill him and his family. Sunil is also a spouse and a father, so private relationship creates certain responsibilities.
\item A moral/ethical dilemma means you have two moral alternatives --- duty to the public vs. responsibility to family. Do not build a one-sided path (e.g., `duty is important, so ignore family'); the answer must address both dimensions.
\end{itemize}
\end{ExampleBox}

\subsection{Dignity in Private Relationships}

Dignity is important in relationships. A husband should not treat his wife as his property, and vice versa; equality matters.

\begin{itemize}
\item Illustration of failure of dignity: a case in Agra where a woman reported her husband missing, cried and made a scene, but had in fact killed her husband and buried him --- an extreme breakdown of the dignity and trust a private relationship demands.
\item Absolutism of dignity: irrespective of background or culture, every human being deserves some dignity. Ultimately, all of us are human beings.
\item Compassion is also essential --- without compassion there is no real relationship; love means at least having compassion.
\end{itemize}

\subsection{Traditions \& Customs and Rule of Law}

\begin{itemize}
\item Traditions and customs govern family relationships --- they tell us what is expected of us (e.g., Raksha Bandhan customs), and we are supposed to follow them.
\item Rule of Law also applies in private relationships. Law governs private relationships too.
\end{itemize}

\begin{GovtPolicy}
\begin{itemize}
\item Laws applicable to private life: Domestic Violence Act, Dowry Prohibition Act (dowry occurs between families), Prevention of Atrocities against SC/ST Act (you cannot use casteist remarks against any person from SC/ST community even if they are your friend).
\item Uniform Civil Code is being brought in various places and will govern private life; the ban on Triple Talaq (which was between spouses, a private relationship) shows law reaching into private life.
\item Our private relationship is governed not just by traditions and customs but also by Rule of Law.
\end{itemize}
\end{GovtPolicy}

\subsection{Commonalities and Differences between Public \& Private Standards}

If we compare the two sets of standards, we find both commonalities and differences.

\paragraph{Commonalities --- and why they exist}
\begin{itemize}
\item Dignity and compassion appear in both; accountability appears in both (from rule of law in private life, and as a public-life value).
\item Reason for commonality: ultimately, whether private or public life, human beings are at the centre. We should treat every human being with dignity.
\item To protect that dignity, there is law: the Domestic Violence Act, Dowry Act, and Prevention of Atrocities Act all protect a person's dignity. Thus dignity, self-respect, individuality and liberty are all guarded by law.
\end{itemize}

\paragraph{Differences --- and why they exist}
\begin{itemize}
\item Traditional customs guide private life but cannot heavily guide public servants. If every place followed its own customs, uniformity of application of law could not be achieved --- so customs should not guide public conduct much.
\end{itemize}

\begin{ComparisonBox}
\begin{itemize}
\item Common to both: dignity, compassion, accountability --- because human beings are at the centre of both, and law exists to protect that dignity.
\item Different: traditions \& customs guide private life but must not heavily guide public servants, as that would defeat uniformity of application of law.
\end{itemize}
\end{ComparisonBox}

\begin{PYQLink}
Anticipated question (this year or next): the importance of ethical standards in private-life relationships --- why ethics in private relationships is all the more important in contemporary times.
\end{PYQLink}

\section{Why Ethics in Private Relationships Matters More Today}

Family and marriage structures are breaking down: spouses are killing each other; children are dumping elderly parents here and there and not discharging their duty. People are becoming more and more individualistic and are not understanding the responsibilities that come with relationships. In this context, ethics in private relationships becomes more important in today's time.

\subsection{Core Reasons}

\begin{itemize}
\item Rising individualism and materialism: relationships are increasingly reduced to transactions --- treating human beings merely as a means. A relationship is maintained only as long as there is self-interest; when it becomes a matter of `business' / self-interest, it becomes difficult, which is why ethics in private relationships is even more relevant today.
\end{itemize}

\begin{ExampleBox}
\begin{itemize}
\item Gleeden app: an app on which people can register to sustain extra-marital relationships, offering anonymity and protection of identity. Reportedly a majority of its users are women, and a large share are from Bengaluru; extra-marital relationships are said to have become much more common and normalized in Bengaluru. The app saw record registrations. Its growth highlights how loyalty and related values are not being followed in relationships --- marriage is losing context.
\item Videos have surfaced of children abandoning or even beating their elderly parents --- evidence of forgetting responsibilities towards family members.
\end{itemize}
\end{ExampleBox}

\subsection{The Concept of Pitr-rin (Debt to Parents / Ancestors)}

Especially in the Gen-Z context, many believe their achievements are entirely the product of their own hard work. When this happens, one becomes ungrateful to the people who played a huge role in one's upbringing --- namely, parents.

\begin{itemize}
\item Pitr-rin is traditionally misunderstood as merely about having a son; it is bigger than that.
\item We tend to think whatever we achieved in life is because of our own hard work, and we discount or take for granted what our parents contributed --- from this arises a lack of the sense of responsibility towards them.
\item As a result, we make decisions keeping only ourselves in mind, forgetting that other people played a huge role in getting us where we are.
\end{itemize}

\begin{KeyConcept}
\begin{itemize}
\item Pitr-rin: the debt we owe to parents/ancestors. Its larger purpose is that we should not live purely individualistically. When we raise a child, much of what we do is done out of sheer love and responsibility, without calculation.
\item The debt to parents can never be fully repaid; but it is honoured by taking pain upon ourselves in doing something for others --- not merely for our own people.
\end{itemize}
\end{KeyConcept}

\subsection{The Fragility of Life --- Lesson on Prioritising Family}

A profound line: you can actually count the number of days you will spend with your sister (or family). Life is very fragile, so `line dropping' --- pausing our ambition to spend time with family --- is very important. We should not postpone the time we have for parents and siblings, because that time is finite.

\begin{itemize}
\item Beyond duty, creating good memories with parents gives huge satisfaction.
\item The least one can do is be honest with one's preparation and one's responsibilities toward parents.
\end{itemize}

\begin{CriticalPoint}
Do not postpone what you should never postpone --- time with parents and siblings. Ambition should not crowd out relationships, because life is fragile and shared time is limited.
\end{CriticalPoint}

\begin{QuickRevision}
Why ethics in private relationships matters more today: breakdown of marriage \& family structures; neglected children and neglected elderly; rising individualism \& materialism; treating humans merely as a means (Gleeden app); loss of the sense of Pitr-rin; and the fragility of life demanding we prioritise finite time with family. When you write the answer, highlight all these dimensions: elderly, husband, wife, children, neglected children, neglected elderly, breakdown of marriage structure.
\end{QuickRevision}

\section{Is a Watertight Separation of Public and Private Life Possible?}

A debate on which a question may be asked. In UP, a public servant having an extra-marital affair was suspended once authorities came to know. Some argue he should be judged only on his public-life conduct, not his private life, and that there should be different ethical standards for public and private life. The core question: Is it possible to have a watertight compartment between public and private life? The answer is NO.

\subsection{Argument 1 --- Virtue Ethics (`We are what we repeatedly do')}

In virtue ethics, whatever you do habitually becomes part of your character. If someone beats his wife regularly in private life, the character being built is one of sheer disrespect for women in general.

\begin{itemize}
\item Persistent unethical behaviour in private life (e.g., domestic violence) will eventually lead to unethical conduct in public life.
\item Consequence: an IPS officer with such a character will not be able to faithfully enforce the Domestic Violence Act --- i.e., poor implementation of the Domestic Violence Act.
\item Real illustration: a senior IPS officer in MP was beating his own wife --- one can imagine his commitment to enforcing the Domestic Violence Act in public life. Ultimately you are the same person in both lives.
\end{itemize}

\subsection{Argument 2 --- Divergence Creates a Crisis of Credibility}

A lot of divergence between public and private life creates a crisis of credibility. Because a public servant is a role model, people look up to and follow them; large deviance between preaching and practice erodes credibility and harms the common good.

\begin{itemize}
\item Keshab Chandra Sen (Bengal social-religious reform movement): got his daughter married very young (at $\sim$13) despite his public advocacy against child marriage --- and thereby lost credibility.
\item Various religious gurus / Babas: accused of doing things in private life that go against what they preach in public --- leading to a crisis of credibility.
\item Gandhiji and the jaggery story: a woman brought her child to Gandhiji because the child ate too much jaggery; Gandhiji asked her to return after two weeks. He first gave up jaggery himself, then advised the child --- because when you preach something, it is important that you purify your own character first.
\end{itemize}

\subsection{Argument 3 --- Higher Standards Required of Civil Servants}

The Code of Conduct for civil servants requires higher ethical standards: one should not do anything, even in private life, that is `unbecoming of a civil servant'.

\begin{itemize}
\item `Unbecoming of a civil servant' means behaviour expected of a civil servant must be maintained; you should not do anything even in your private life that is unbecoming of the office.
\item Positive model --- Lal Bahadur Shastri: when the country faced difficulty (drought, etc.), he first started skipping a meal himself in his own private life before asking others to fast --- showing that people follow only when leaders start with their own conduct.
\item Additional illustrations: AP suspended an IPS officer for unauthorised foreign visits; the MP IPS case; and `imagine \rupee 200 crore for dowry' --- what example does that set for the youth? A private affair can have huge consequences on public life.
\end{itemize}

\begin{KeyConcept}
Watertight separation between public and private life is NOT possible, because (1) virtue ethics --- we are what we repeatedly do; (2) large divergence creates a crisis of credibility; and (3) civil servants are held to higher standards and must avoid conduct `unbecoming' of the office.
\end{KeyConcept}

\begin{PYQLink}
\begin{itemize}
\item 2023: `Corruption is the manifestation of the failure of core values in the society.' --- you can never eradicate corruption without solving the problem of values in society.
\item Related theme: whether public servants should be judged on private conduct --- build the answer using virtue ethics, crisis of credibility, and the code of conduct.
\end{itemize}
\end{PYQLink}

\subsection{Where to Draw the Line: Do Public Officials Have Any Privacy?}

If we cannot separate the two, does it mean civil servants have no right to privacy --- that every aspect of private life should be scrutinised, especially in the age of social media? The answer is NO --- private life does exist. The line is drawn by impact on public duty.

\paragraph{Two contrasting illustrations}
\begin{itemize}
\item Not to be scrutinised: if posted where alcohol is not illegal, having a drink with a friend after office is private conduct that does not affect duty.
\item To be scrutinised: drinking all night and reaching office hungover --- because it now impairs the discharge of public duty.
\end{itemize}

\begin{CriticalPoint}
\begin{itemize}
\item The principle: to the extent private life influences the discharge of public duty, it should be scrutinised.
\item Examples that MUST be scrutinised because they bear on public responsibility: whether an official is taking dowry; gambling to the point of bankruptcy (which increases proneness to corruption). Conduct with no bearing on duty remains genuinely private.
\end{itemize}
\end{CriticalPoint}

\begin{QuickRevision}
Two big exam themes: (1) Why ethics is required in private relationships --- given rising individualism, materialism and family breakdown. (2) Whether a watertight separation exists between public \& private life (No), and if not, whether any privacy remains (Yes --- only conduct that influences public duty is legitimately scrutinised).
\end{QuickRevision}

\section{Human Values}

Earlier we discussed certain human actions and asked whether they qualify as `human' actions. Acts such as manual scavenging, bonded labour, one person urinating on another (an MP incident), and slavery are inhuman actions --- because when interacting with the other person, we are not considering them as a human being.

\subsection{Meaning of Human Values}

\begin{KeyConcept}
\begin{itemize}
\item Human values are those values that guide people to take into account the human element --- the fact that the other person is human --- when interacting with each other.
\item Human dignity is a very important, absolute value.
\end{itemize}
\end{KeyConcept}

\subsection{Controlling Emotions and Desires --- Plato's Chariot Analogy}

Controlling your desires is a human value. If a person is driven by desires, they behave more like an animal. Plato's analogy: you (the charioteer / soul) are drawing a chariot with two horses --- one horse is emotions and the other is desires; both drive you.

\begin{itemize}
\item When desires drive you: people commit rape and indulge in greed.
\item When emotions drive you: emotions like anger lead to mob lynching, honour killing, or throwing acid on a girl when a proposal is rejected.
\item As human beings, it is expected of us that we control both horses --- emotions and desires. Losing control is a loss of humanity.
\item We should be fair to each other, because all of us are human beings.
\end{itemize}

\begin{CriticalPoint}
Life is a long journey between `being human' and `being humane'. Not every human action reflects humanity --- humanity is about what we ought to be.
\end{CriticalPoint}

\subsection{Types of Values: Terminal vs Instrumental}

Values are of two types.

\begin{center}
\begin{tabularx}{\textwidth}{>{\raggedright\arraybackslash}p{2.9cm} Y Y}
\rowcolor{AccentDark}
\thead{Aspect} & \thead{Terminal Values} & \thead{Instrumental Values} \\
\toprule
Meaning & End-goal values (the `ends') & Means / values that are your chosen ways of behaving \\
Examples & Inner peace, happiness, equality, family security, national security, freedom, goals & Responsibility, self-control, honesty, tolerance \\
Role & The destination we are pursuing & The behaviour we adopt to reach that destination \\
\bottomrule
\end{tabularx}
\end{center}

\begin{itemize}
\item Terminal values are the ends --- e.g., inner peace, happiness, equality, family security, national security, freedom/goals, tolerance of different perspectives.
\item Instrumental values are the means / chosen behaviours --- e.g., doing our work with responsibility, exercising self-control.
\end{itemize}

\begin{QuickRevision}
Human values take the human element into account when we interact. Human dignity is an absolute value. Controlling emotions \& desires (Plato's chariot) is central. Values split into terminal (ends: inner peace, happiness, equality, security, freedom) and instrumental (means: responsibility, self-control, honesty, tolerance).
\end{QuickRevision}

\section{Inculcation of Values: Why Values Must Be Inculcated}

Any society wants us to behave well and conduct ourselves rightly, and for that it makes laws --- government makes laws. But can good behaviour be ensured by law alone? No, it cannot.

\begin{itemize}
\item Despite laws, corruption still happens; despite laws, dowry still happens. Law alone is insufficient.
\end{itemize}

\subsection{The Tanker-Looting Illustration}

\begin{ExampleBox}
\begin{itemize}
\item After an accident, people looted the overturned tanker. The government did not ask anyone to loot it --- yet people did. We usually blame corruption on the government and forget that we (society) also have a role.
\item Personal account: after the speaker's own accident, the ambulance attendant who dropped him at the hospital asked whether he could take the goods from the wrecked car; and when the car was later recovered from the accident site, the speaker, the extra tyre, and other fittings had been stolen --- even though everyone knew it was an accident-damaged vehicle.
\end{itemize}
\end{ExampleBox}

The lesson: most people are against corruption not out of principle but because they lack the opportunity; the moment they get the chance, they may do it. Humans can behave like `vultures'. If we want a country without corruption, it must come from society --- law alone will not help.

\begin{itemize}
\item P. Sainath, `Everybody Loves a Good Drought' --- the title captures how, during a drought, relief money brings joy because relief money reaches private pockets; corruption is also a societal problem, not just a government problem.
\item Lokpal debate (UPSC asked): a national Lokpal / stronger bill alone will not solve corruption if it is collusive corruption --- where those with money to bribe escape. The root cause is societal values.
\end{itemize}

\subsection{Education Without Values $\rightarrow$ `A Clever Devil'}

Government can make good laws and provide documents, but ultimately we have to use them well. Education provided without values will make a person a `more clever devil'.

\begin{CriticalPoint}
Education without values makes a person a more clever devil: they will misuse acquired skills, and the consequences of that misuse are far greater than if they had not been educated at all (e.g., a doctor who becomes a drug prescriber; skilled people launching cyber-attacks).
\end{CriticalPoint}

\begin{itemize}
\item On any topic --- women, corruption, etc. --- your suggestions will not be complete without moral dimensions; behind all these problems lies a problem of morality.
\item On demand-supply and corruption: even if demand is high and supply is low, that does not entitle anyone to do corruption. Corruption exists because we lower our moral guard --- there is always a choice (you can earn food through effort rather than by snatching).
\end{itemize}

\subsection{Many Problems Are Societal, Not Merely Legal}

\begin{itemize}
\item Conductor assaulted for not speaking in Marathi --- a societal problem, not only a legal one.
\item Fake cardiologist booked over patient deaths --- someone falsely posing as a heart specialist and even operating on hearts; the audacity shows the level of moral failure (societal problem).
\item Constitutional morality is not a natural sentiment --- it has to be cultivated, precisely because our socialised instincts are not automatically constitutional.
\end{itemize}

\begin{PYQLink}
2023: `Corruption is the manifestation of the failure of core values in the society.' You will never be able to eradicate corruption without solving the problem of values in society.
\end{PYQLink}

\begin{QuickRevision}
Why inculcate values: law alone cannot ensure good conduct (corruption, dowry persist despite law; the tanker-looting shows society's own role). Most oppose corruption only for lack of opportunity. Education without values = a clever devil. Many problems (Marathi-assault, fake cardiologist, corruption) are societal, not merely legal. Constitutional morality must be cultivated.
\end{QuickRevision}

\section{Role of Family in Value Inculcation}

Family, education institutions, and society all play a huge role in inculcating values. We begin with the family, illustrated through great examples that are useful in both answers and essays.

\subsection{Examples of Values Inculcated by Family}

\begin{itemize}
\item Ratan Tata --- Altruism \& Philanthropy: his altruism and philanthropy came from the Tata family, which historically had a legacy of philanthropy and helping people.
\item Nelson Mandela --- Forgiveness \& Reconciliation: known for forgiveness. Belonging to a political family, he saw people come daily to the local/village council (panchayat) to resolve their differences through discussion and reconciliation. This early exposure to village councils taught him to build society by reconciliation, not revenge. He said that if he had not forgiven, he would still be in prison --- because holding onto the past keeps you imprisoned in it.
\item B.\,R. Ambedkar --- Education as empowerment: his father, a Subedar in the army, insisted on educating his children despite social rules against Dalits getting educated. Ambedkar learnt the importance of education as a tool of empowerment --- reflected later in the architecture of the Constitution.
\item A.\,P.\,J. Abdul Kalam --- Simplicity \& Tolerance: learnt simplicity from his boatman father, and tolerance from teachers of different faiths who welcomed him (a Muslim) --- which is why he later became the `People's President'.
\item Martin Luther King Jr. --- Speaking against racism: his father was a pastor/religious authority who spoke against racism in his sermons; from there, King's value of fighting against racism was instilled.
\item Chhatrapati Shivaji Maharaj --- role of Jijabai: Jijabai's role shaped his rise and values.
\item Lord Ram --- Upholding promises: learnt to uphold promises from his family; his father honoured his word even at great cost, and Ram fulfilled his promise even at the cost of exile.
\end{itemize}

\subsection{Values from Joint Families}

\begin{itemize}
\item A joint family inculcates cooperation, religious values, responsibility, compassion, discipline, and respecting elders.
\end{itemize}

\begin{QuickRevision}
Family examples (keep 6--8 ready for answers/essays): Ratan Tata (altruism/philanthropy), Nelson Mandela (forgiveness/reconciliation via village councils), B.\,R. Ambedkar (education as empowerment), A.\,P.\,J. Abdul Kalam (simplicity \& tolerance), M.\,L. King Jr. (against racism), Shivaji \& Jijabai, Lord Ram (upholding promises). Joint families add cooperation, discipline, responsibility, respecting elders.
\end{QuickRevision}

\subsection{How Families Inculcate Values (Mechanisms)}

\paragraph{Child-rearing practices} How a child is reared has huge consequences for character development.

\begin{ExampleBox}
\begin{itemize}
\item Ruth Benedict --- `The Chrysanthemum and the Sword': Americans were astonished that Japanese soldiers were disciplined yet ruthless. Benedict studied Japanese culture: a newborn is treated like a flower (chrysanthemum) with tenderness and care, but as the child grows, heavy discipline is imposed --- even toiletry practices are regulated, and emotional expression is suppressed. The child becomes disciplined but ruthless, because their own emotions were never understood.
\item Serial killers often become so because of a traumatic past --- trauma from childhood scars the character. Contrast: authoritarian vs democratic upbringing.
\end{itemize}
\end{ExampleBox}

\paragraph{Other mechanisms}
\begin{itemize}
\item Observation learning: a child learns by watching --- e.g., patriarchy and smoking are learnt by observation.
\item Storytelling: stories with a moral lesson --- Ramayana, Mahabharata, Panchatantra, the rabbit and the tortoise.
\item Operant conditioning: reward--punishment mechanism. Good behaviour (hard work, good marks) is rewarded, teaching desirability, which becomes a preference and then a value; wrong acts (stealing, lying) are punished, teaching what not to do.
\item Role modelling: children learn a lot from parents, who become their role models.
\end{itemize}

\subsection{Critical Perspective --- Why Families Fail to Inculcate Values (Contemporary Factors)}

\begin{itemize}
\item Parents are not great role models: they indulge in unethical acts (e.g., jumping traffic signals) in front of their children (observation learning of the wrong kind).
\item Breakdown of the family / joint-family system. Parenting is a selfless duty, but at times parents take career to such a level that this is compromised.
\item Fast-paced life and working parents who may not be emotionally available for the child.
\item Social media \& online games: at a very young age children are exposed to inappropriate content (some countries have even banned social media for under-16s). Even `good' content on social media does not develop emotional intelligence, because you are not interacting with a real human being (missing non-verbal communication). It also reduces human interaction. Examples: the Blue Whale game; the Netflix series `Adolescence'.
\item Parents undermining ethics directly: e.g., paying for leaked exam papers --- emphasising `marks over morality' and `career over character'.
\end{itemize}

\begin{CriticalPoint}
\begin{itemize}
\item Cognitive atrophy from phones: parents hand children phones to reduce their own burden; $\sim$90\% of the adult-size brain develops in the first two years, and heavy phone exposure harms this development --- a growing crisis in cognitive ability.
\item Key phrases for answers: `Marks over morality', `Career over character', emphasis on career over character.
\end{itemize}
\end{CriticalPoint}

\begin{PYQLink}
UPSC (2017 \& 2022-type, near copy-paste): a question on the role of family (parents) and teachers in value inculcation. A good answer must include contemporary factors to make it analytical, not merely positive.
\end{PYQLink}

\section{Role of Educational Institutions in Value Inculcation}

\subsection{Role of Teachers}

Teachers play a huge role. Traditionally many teachers focused on rote learning (`rattafication'), which limited genuine understanding. A conceptually-oriented teacher who explains why an operation is invalid --- rather than only giving the solution --- transforms students. Such teachers act as agents of transformation.

\begin{itemize}
\item Pedagogy is very important: the subject is not merely about transferring information; students should `live' and feel the subject. Teachers shape the character of the student through their teaching, behaviour, and conduct.
\item Professor H.\,C. Verma (awarded Padma Shri): his book `Concepts of Physics' helped countless students understand complex concepts and pursue science careers; he inculcated scientific temper and inspired dream careers through role modelling.
\item Anne Sullivan \& Helen Keller: Sullivan taught a child with hearing and visual impairment --- showing the empathy and pedagogy required to reach such a student. Teachers stimulate critical thinking, objectivity and other values, and mentor students throughout their journey.
\end{itemize}

\begin{KeyConcept}
Teaching carries responsibility far beyond examinations: through pedagogy, behaviour and role modelling, teachers inculcate values and can transform a student's entire perspective and life.
\end{KeyConcept}

\begin{itemize}
\item More teacher--mentor examples: Ramakrishna Paramahamsa \& Swami Vivekananda; Samarth Guru Ramdas \& Chhatrapati Shivaji Maharaj; Chanakya (Kautilya) \& Chandragupta Maurya.
\end{itemize}

\subsection{Role of Curricular Subjects}

\paragraph{History --- lessons, not just facts}
\begin{itemize}
\item Justice cannot be kept by force: force may create an absence of tension but not genuine, long-lasting peace. The Treaty of Versailles, imposed by force rather than justice, helped create the Second World War.
\item Tolerance --- Akbar vs Aurangzeb: Akbar forged a large empire through a tolerant approach (e.g., Din-i-Ilahi). Aurangzeb, with an intolerant approach --- discriminating between Shias and Sunnis and imposing Jizya --- could not forge a lasting empire. Under intolerance an empire may survive out of fear, but people do not connect from the heart; after his death, the Mughal empire crumbled.
\end{itemize}

\paragraph{Science --- explaining, questioning, predicting}
\begin{itemize}
\item Science explains phenomena (tides, eclipses) that were once attributed to Rahu and Ketu; it not only explains but predicts them --- enabling cyclone management in Odisha and drastically reducing deaths.
\item Superstition example: a snake-bitten child in UP was tied and floated down the Ganga instead of being taken for treatment, and died --- ignorance because the phenomenon was not understood scientifically.
\item Scientific temper / critical thinking theme (very important this year): `A doubter is a true soldier of science' (the doubter questions); `Visionary decision-making happens at the intersection of intuition and logic' (logic = questioning); `History is a series of victories of the scientific man over the romantic man'; `Thinking is like a game that does not begin unless there is an opposition.'
\item Personal chilli illustration: an itching sensation from chilli oil, relieved by soap (fat/carbon counteracting the oil) --- science makes us look into and understand things rather than take them at face value.
\end{itemize}

\paragraph{Environment}
\begin{itemize}
\item Studying environment --- climate change, global warming, groundwater depletion --- teaches the importance of environmental conservation.
\end{itemize}

\subsection{Extracurricular Activities: Sports, Games \& Debates}

Sports and games are a very important --- and under-used --- source of ethics, so they are worth quoting for a distinct answer or essay.

\begin{itemize}
\item Facing failure: we cannot win all the time; a key life skill is how we face failure. Sports teach us to digest failures and bounce back, which builds character.
\item Sportsmanship: respecting and appreciating others' tremendous effort, even in defeat.
\item Teamwork \& selflessness: putting the team's interest above individual ego --- a student who, as captain, set aside individual defences for the team learnt selflessness (linked to public duty). Graeme Smith batting with a broken hand in a Test to save his team is an instance of putting the team above oneself.
\item Administration using sport: the Mewat Kabaddi league was used to change attitudes and fight drug addiction --- sport can influence people's lives.
\item Debates: engaging in debate makes you more tolerant and accepting of others' opinions and respectful of liberty --- you realise your opinion may not be the best.
\end{itemize}

\begin{PYQLink}
Essay theme (recent): `The most valuable lessons in life are learned through bitter experiences.' Use sport, e.g., the 2023 World Cup final loss --- allowing expectations to rule our head and trying to control what is not in our control means we may not give our best.
\end{PYQLink}

\subsection{Pedagogy: Activity-Based Learning \& the Flipped Classroom}

\begin{itemize}
\item Activity-based learning (observed in Norway): students solve problems together through projects --- a higher form of learning than lecture-only teaching.
\item Flipped classroom: in a traditional classroom, lectures come first and homework after; in a flipped classroom, you study a topic on your own (e.g., recorded video) beforehand and then solve problems together through activity in class.
\end{itemize}

\subsection{Critical Perspective --- Why Institutions Fail (Contemporary Factors)}

\begin{itemize}
\item Unethical conduct by teachers/institutions: e.g., discrimination by caste --- making students sit separately during the mid-day meal --- inculcates inequality; sexual abuse; and other unethical acts.
\item Commercialisation, commodification \& privatisation of education.
\item Over-emphasis on marks and exam performance.
\item Neglect of extracurricular activities --- missing out on the values sports and games inculcate.
\item Commercialisation of curriculum; neglect of the psychological well-being of students; students indulging in bullying, slang and ragging.
\item Lack of teacher training / pedagogical bankruptcy --- emphasis on rote learning.
\end{itemize}

\begin{GovtPolicy}
\begin{itemize}
\item Positive institutional interventions (usable examples): the Malayalam movie `Sthanarthi Sreekuttan'-inspired semi-circular seating layout in Kerala schools to avoid labelling of `backbenchers'; the sugar board displaying sugar content to make students health-conscious; NCERT emphasis on traditional games and cultural activities (kabaddi, yoga) to impart teamwork, agility and other values.
\item Activity/project-based learning gap: most Indian institutions are `teaching institutions' relying on lectures; project- and activity-based learning could impart more values and knowledge --- a missed opportunity, alongside the prominent neglect of extracurricular activities.
\end{itemize}
\end{GovtPolicy}

\begin{QuickRevision}
Educational institutions inculcate values via teachers (pedagogy, role modelling --- H.\,C. Verma, Anne Sullivan), curricular subjects (History $\rightarrow$ tolerance/justice: Versailles, Akbar vs Aurangzeb; Science $\rightarrow$ question--explain--predict, scientific temper; Environment $\rightarrow$ conservation), and extracurriculars (sports $\rightarrow$ failure, sportsmanship, teamwork, selflessness; debates $\rightarrow$ tolerance). Better pedagogy = activity-based learning \& flipped classroom. Failures: caste discrimination, commercialisation, marks-obsession, neglect of extracurriculars, weak teacher training.
\end{QuickRevision}

\section{Role of Society, Government \& Media in Value Inculcation}

Value inculcation in society is contingent not just on family and educational institutions; government, media, and religious leaders all play a huge role.

\subsection{Ethnocentrism \& the Value of Exposure (Society)}

\begin{ExampleBox}
\begin{itemize}
\item The `frog-in-the-well' realisation: raised in a traditional vegetarian Brahmin family in Rajasthan, the speaker once judged non-vegetarians as `wrong people'. At IIT Kharagpur, through a seniors' `personality program' (remembering each person's name, department and place), he came to know roommates and messmates from Bengal (fond of fish), Hyderabad (mutton), and Meghalaya. He realised that had he been born in Calcutta, Hyderabad or Meghalaya, he too would eat those foods.
\item Insight: dietary practices are shaped by the culture into which we are born --- a culture we do not choose. Without exposure to diversity, we become ethnocentric and judge others by our own standards. This is why `morality is not a natural sentiment' --- it is shaped by socialisation. Educational institutions are great places to enlarge our thinking.
\end{itemize}
\end{ExampleBox}

\subsection{Empathy Toward Systems and People}

When operating at a small scale you can ensure things are done well; but as part of a larger nation/system you realise it is not humanly possible to do everything perfectly. Hence we should extend empathy --- assume good intentions and recognise limitations --- rather than only understanding our own problems and not the other's perspective.

\subsection{Role of Government}

Government mostly inculcates values through a reward--punishment approach, plus attitudinal change and awareness.

\begin{itemize}
\item Rewarding ethical conduct: e.g., civilian awards signal that a behaviour is appreciable.
\item Punishing unethical conduct through laws: e.g., Domestic Violence Act, anti-corruption law.
\item Striving for attitudinal change and building awareness.
\end{itemize}

\begin{GovtPolicy}
Incentivising behaviour: government announcing a \rupee 1 crore incentive for a `Maoist-free panchayat' --- using incentives to shape conduct.
\end{GovtPolicy}

\subsection{Role of Media (Positive)}

\begin{itemize}
\item Shaping public opinion \& spreading awareness: e.g., coverage of the Mooknayak-type / `money-per' incident sensitises us (we did not go to Manipur, yet media sensitised us --- `what the eyes do not see, the heart does not grieve about').
\item Solidarity \& cooperation: e.g., appeals to donate during Punjab floods.
\item Reinforcement of shared values \& traditions: e.g., 26th January parade celebrating unity in diversity.
\item Promoting tolerance: e.g., debates and discussions; popularisation of South Indian movies in North India.
\item Inculcating scientific temper: e.g., media advisories to farmers for scientific agriculture.
\item Eradicating stereotypes: e.g., the movie `Pink' --- `No means no'.
\item Upholding justice \& accountability: e.g., coverage of scams (2G), the NEET / paper-leak issue --- holding government accountable.
\end{itemize}

\subsection{Role of Media (Negative --- Contemporary Factors)}

\begin{itemize}
\item Objectification \& commodification of women: how women are portrayed in songs and movies shapes understanding (e.g., a Haryanvi song for which the singer later apologised).
\item Echo-chamber effect: strengthening one's existing prejudices; fake news gets reinforced when you lack alternative perspectives.
\item Biased media coverage \& the algorithm problem --- undermining objectivity and critical thinking.
\item Prioritising sensationalism over public interest: pursuing sensational stories over the public interest. What is popular may not be right --- popularity is not the metric of righteousness.
\item Loss of media autonomy due to corporate takeover; glorification of crime and criminal syndicates --- glorifying something that goes against rule of law.
\end{itemize}

\begin{QuickRevision}
Society enlarges thinking and counters ethnocentrism through exposure (frog-in-the-well). Government inculcates values via rewards (civilian awards), punishment (laws), and awareness/incentives (\rupee 1 crore for Maoist-free panchayat). Media positives: sensitising, solidarity, shared values (Republic Day), tolerance, scientific temper, breaking stereotypes (`Pink'), accountability (2G, NEET). Media negatives: objectification of women, echo chambers, bias/algorithms, sensationalism, glorification of crime, loss of autonomy.
\end{QuickRevision}

\section{Lessons from the Lives of Great Leaders --- Mahatma Gandhi}

\subsection{Honesty \& Integrity}

\begin{itemize}
\item Paying for his grandchildren's stay: when Gandhiji's grandchildren came to stay with him in the ashram, he deposited fees for their stay because the ashram was run from public funds --- a great example of honesty and integrity.
\item Expelling his sister over untouchability: Gandhiji invited Dudabhai (a Harijan teacher) to live in the ashram; when his sister refused to eat with Dudabhai, Gandhiji acted against the discriminatory refusal --- a great example of honesty and integrity in upholding equality over custom.
\end{itemize}

\subsection{Openness, Tolerance \& a Learning Attitude}

Gandhiji used to say we should keep our mind open like we keep the windows of our house open, so that sunlight and wind can come in --- similarly, we should keep our mind open so we can learn from others.

\begin{KeyConcept}
\begin{itemize}
\item Gandhiji said he learned non-violence (Ahimsa) from his wife: he had always tried to make her bend to his wishes; she would firmly refuse yet patiently bear the hardship he inflicted. Her peaceful opposition opened his eyes, made him ashamed, and rid him of the notion that it was his birthright to rule over her.
\item Lesson --- openness, tolerance and a learning attitude: even a great person admits his mistakes; great people own them.
\end{itemize}
\end{KeyConcept}

\subsection{Non-Violence (Ahimsa)}

\begin{itemize}
\item Non-violence is not only refraining from harming others through action, but also through our thoughts and words. Use this whenever you criticise hate speech, abuse or violence in case studies.
\item Withdrawal of the Non-Cooperation Movement after the Chauri Chaura incident --- a powerful example of Gandhiji's commitment to non-violence.
\end{itemize}

\subsection{Forgiveness}

Forgiveness is a big enough theme to become a standalone question --- and whenever you write on forgiveness, always name two persons: Gandhiji and Nelson Mandela.

\begin{itemize}
\item `Hate the sin, not the sinner' did not mean people should not be punished --- Gandhiji even punished himself (fasting) when he felt he had done wrong. He was not against punishment; he was against bitterness. Punishment should be for rehabilitation and reform, not out of hatred (just as parents scold to correct, hating the mistake and not the child).
\item Mandela: `If I had not forgiven them, I would still be in prison.' If we do not forgive, we remain imprisoned in grudges and keep reliving the wrong. Forgiveness liberates the one who forgives --- you forgive first for yourself.
\end{itemize}

\subsection{Swaraj (Self-Rule)}

\begin{itemize}
\item Swaraj literally means self-rule / rule over oneself. The greatest self-rule is ruling over oneself --- regulating oneself. Mention Swaraj whenever you write on emotional intelligence or on controlling desires.
\end{itemize}

\subsection{Talisman \& Antyodaya}

\begin{KeyConcept}
\begin{itemize}
\item Gandhiji's Talisman: whenever you are in doubt during decision-making, think of the poorest and most vulnerable (`last-mile') person and ask how your action affects them --- take into account the most vulnerable (Antyodaya).
\item Antyodaya = feeling others' pain AND doing something to remove it. Example: during the lockdown, government did not sufficiently anticipate the hardship of migrants --- the Talisman perspective would have centred the most vulnerable.
\end{itemize}
\end{KeyConcept}

\subsection{Solidarity --- The Loincloth}

\begin{itemize}
\item Adopting the loincloth: originally a barrister who wore full clothes, Gandhiji, while travelling by train in South India, urged people wearing foreign clothes to switch to khadi; they replied that khadi was too costly and they could not afford it. Realising the poor could not afford khadi, and to show solidarity with their pain, he changed his own lifestyle and began wearing the loincloth --- a profound act of solidarity, feeling and sharing the pain of others.
\end{itemize}

\subsection{Sarvodaya, Purity of Means \& Trusteeship}

\begin{itemize}
\item Sarvodaya = upliftment / rise of all (`sabhi ka uday'). Gandhiji spoke for everyone --- farmers, Dalits, women --- linked to inclusive growth.
\item Purity of means: the means must be pure --- illustrated by the withdrawal of the Non-Cooperation Movement after Chauri Chaura.
\end{itemize}

\begin{KeyConcept}
\begin{itemize}
\item Trusteeship: if you have resources beyond your genuine need (e.g., you need \rupee 100 to survive but have \rupee 1,000), the extra is not yours --- you hold it as a trustee on society's trust.
\item Application --- Bhoodan movement: after independence, most land lay with zamindars/landlords; Vinoba Bhave appealed to those with excess land to donate it. Bhoodan reflects the trusteeship principle.
\item Relevance today: trusteeship matters amid rising inequality (the top 1\% holding vast wealth) and unsustainability --- regarding the environment, we hold all resources as trustees of future generations, so trusteeship underpins environmental conservation. (Also links to CSR.)
\end{itemize}
\end{KeyConcept}

\begin{QuickRevision}
Gandhian values for answers/essays: Honesty \& Integrity (ashram fees; expelling sister over untouchability); Openness/Tolerance/Learning (mind like open windows; learnt Ahimsa from his wife); Non-Violence (withdrawal of NCM after Chauri Chaura); Forgiveness (`hate the sin not the sinner'; with Mandela); Swaraj (self-rule $\rightarrow$ emotional intelligence); Talisman \& Antyodaya (last-mile person); Solidarity (loincloth); Sarvodaya (upliftment of all); Purity of Means; Trusteeship (Bhoodan; inequality; environment; CSR).
\end{QuickRevision}

\section{Key Cross-Cutting Reminders}

\begin{itemize}
\item Not all values are ethical values. Ethics is about standards; values are broader. (Code of ethics vs code of conduct will be covered separately.)
\item Case studies \& dilemmas: when public duty and private responsibility clash, resolve via the golden mean that addresses both moral alternatives, not one side alone.
\item Make answers analytical: on family/education/society/government/media, always pair positive roles with contemporary factors (why values are not being inculcated) to lift the answer beyond a purely positive account.
\end{itemize}

\begin{PYQLink}
High-probability themes to prepare: (1) role of family, teachers \& society in value inculcation (UPSC 2017/2022); (2) `Corruption is the manifestation of the failure of core values in society' (2023); (3) `Most valuable lessons are learned through bitter experiences' (essay); (4) forgiveness as a standalone value (Gandhi + Mandela); (5) ethics in private relationships \& the limits of scrutinising a public servant's private life.
\end{PYQLink}
